\section{Full-boundary assembly of compressed surface covers}
\label{sec:full-boundary-assembly}

The maintained \texttt{surface\_cover.py} classifies a supplied dihedral
cover of one canonical bordered surface.  Its component and marking formulas
do not themselves specify what happens when entire collections of boundary
circles are attached.  The new \texttt{surface\_gluing.py} implements that
operation under an explicit promise: every connected collection of patches
has one common sheet coordinate, every monodromy and seam transport is
dihedral in that coordinate, and each seam glues a complete base boundary
together with its full inverse image.  The assembled base may be closed.

The mathematical ingredients are classical covering-space classification,
van Kampen's theorem, and the dihedral arithmetic already used by the
repository.  The contribution is a precise closure theorem for this checked
input model, a provenance-preserving implementation, and independent tests of
its topology and marked transport.  It is not a new general orbit-counting
theorem: Agol--Hass--Thurston's more general compressed normal-surface orbit
algorithms already have polynomial bit complexity.  Nor does this kernel
produce an unknot verdict or assert that an assembled surface is embedded in
a specified three-manifold.

\subsection{Input, boundary directions, and the seam equation}

Let $X=\mathbb Z/W\mathbb Z$, where $W\ge1$ is given in binary.  For every
vertex $v$ of a finite graph, let $F_v$ be a connected compact surface with
at least one boundary circle.  Its orientability, genus or crosscap number,
and number of boundaries specify the same canonical free generators and
peripheral words as the existing cover interface.  Assign to every generator
an affine permutation
\[
 g_{v,i}(x)=\sigma_{v,i}x+a_{v,i},\qquad \sigma_{v,i}\in\{1,-1\},
\]
and let $w_{v,i}\in\mathbb Z/2$ be the actual orientation character of that
generator.  The bit $w_{v,i}$ is determined by the surface schema; it is
not inferred from $\sigma_{v,i}$.  Write $B_{v,j}$ for the evaluated
monodromy of its $j$th canonical peripheral word.

Each patch fixes a local orientation at its basepoint, a marked reference
point on each boundary, and a whisker from the basepoint to that point.
The boundary direction is induced by the collar orientation transported
along this whisker.  The usual polygon models with handle commutators or
crosscap squares realize these conventions.  Written words are traversed
in their written order; affine expressions below use functional composition.

A seam $e$ specifies two distinct unused boundary ports $(u,a)$ and $(v,b)$,
a boundary degree $\epsilon_e\in\{1,-1\}$, and a bijection of reference
fibres
\[
 A_e(x)=s_ex+c_e,\qquad s_e\in\{1,-1\}.
\]
It is admitted precisely when
\begin{equation}
 A_e\circ B_{u,a}= B_{v,b}^{\epsilon_e}\circ A_e.
 \tag{S}\label{eq:seam-intertwining}
\end{equation}
Every boundary port occurs in at most one seam.  A self-seam is allowed
only between two different boundaries of the same patch.  No selected lift
circle or proper interval of a boundary is silently substituted for a full
inverse image.

Equation~\eqref{eq:seam-intertwining} is equality of permutations of $X$,
not equality of formal sign--shift pairs.  Two affine expressions agree
exactly when both their coefficient difference and their constant difference
vanish modulo $W$.  This distinction is essential for $W=1,2$, where formal
signs $+1$ and $-1$ can describe the same permutation.

\paragraph{Realization of one seam.}
A covering of a circle is the mapping torus of its monodromy permutation.
A fibre bijection $A$ extends across a degree-$\epsilon$ map of the base
circle if and only if it intertwines the monodromy as in (S).  For example,
in a parameter interval one lifts the prescribed base homeomorphism from
the image of every reference-fibre point; (S) is exactly the condition that
the endpoints match on identifying the ends of the interval.  This produces
a homeomorphism between the full boundary covers.  Its inverse is obtained
from $A^{-1}$ and the inverse base map.  Thus every accepted seam is an
actual surface attachment, including when its lift circles have large
covering degree.  After gluing all disjoint boundary pairs, collars give a
compact surface and the local covering charts join to a $W$-sheeted cover
of the sewn base.

\subsection{A spanning-tree gauge and the exact assembled monodromy}

Split the patch adjacency graph into connected components.  The following
construction is performed on each component separately, with root $o$ and
spanning tree $T$.  Let $P_v:X_o\to X_v$ be the affine transport along
the tree from $o$ to $v$, using $A_e^{-1}$ when traversing a seam backwards.
Let
\[
 \delta_e=\begin{cases}0,&\epsilon_e=-1,\\1,&\epsilon_e=+1,\end{cases}
 \qquad
 t_v=\sum_{e\in[o,v]_T}\delta_e\pmod2.
\]
The bit $\delta_e$ follows directly from the boundary-orientation convention:
orientation frames extend consistently across a seam that reverses the
induced boundary directions.  The fibre sign $s_e$ has no role in this test.

Use the following based permutations at the root:
\begin{align*}
 \widehat g_{v,i}&=P_v^{-1}\circ g_{v,i}\circ P_v,
       &\widehat w_{v,i}&=w_{v,i},\\
 H_e&=P_v^{-1}\circ A_e\circ P_u\quad(e\notin T),
       &\eta_e&=t_u+\delta_e+t_v\pmod2.
\end{align*}
For every unglued boundary $(v,b)$ retain
\[
 \widehat B_{v,b}=P_v^{-1}\circ B_{v,b}\circ P_v.
\]
Conjugation and multiplication preserve the dihedral affine form.  The
number of root generators is the sum of the local free ranks plus the
cycle rank of the patch graph.  They can be redundant; no freeness assertion
is made for a closed assembled base.

\begin{proposition}[Root action of a sewn cover]
The connected components of the assembled cover are exactly the orbits
of the group generated by the $\widehat g_{v,i}$ and $H_e$.  Its orientation
character on these loop generators is respectively $w_{v,i}$ and $\eta_e$.
The permutations $\widehat B_{v,b}$ are the peripheral monodromies of its
remaining boundaries.
\end{proposition}
\begin{proof}
The local spines, their boundary whiskers, and the seam-crossing paths
generate the fundamental group of the sewn surface by van Kampen.  Tree
paths change the basepoint and yield the displayed conjugations; non-tree
edges supply the remaining graph loops.  The relations coming from a glued
circle hold because (S) was checked before assembly.  The induced covering
therefore has precisely the displayed monodromy action.  Standard path
lifting identifies connected covering components with its orbits.

Alternatively, form the explicit graph with vertices $(v,x)$, local
generator edges, and seam edges $(u,x)\leftrightarrow(v,A_ex)$.  Transport
along a tree identifies every vertex with a root sheet, and leaves exactly
the displayed local and non-tree root actions.  This proves the orbit claim
without constructing a presentation for the closed surface.  Tracking a
local orientation frame along the same paths gives the character formulas;
the two copies of the tree path cancel around a local generator.  The
peripheral formula is the same change-of-basepoint calculation.
\end{proof}

The code stores the original patch and seam data together with the gauge
forest.  A separate verifier checks parent incidence, acyclicity, root
coverage, the affine transport equations, and the orientation parity of
each tree edge.  It checks every parent chain once.  The verifier permits
another valid tree and another root; it never calls the producer's forest
builder or the dihedral classifier.  This is a certificate of tree
transport, not an independent certificate of arbitrary serialized genus
and component claims.  Those claims are recomputed from the input.

\subsection{Component topology, including closed output}

Let $G$ be the resulting affine group at one root.  If $G$ has no supplied
reflection, put $d=\gcd(W,a_1,\ldots,a_r)$.  Otherwise choose one reflection
$R(x)=-x+a_0$ and take the gcd of $W$, all translation shifts, and all
differences of reflection shifts from $a_0$.  Put $m=W/d$.  Products of
two reflections and Bezout's identity give
\[
 G=\langle x\mapsto x+d\rangle
 \quad\hbox{or}\quad
 G=\langle x\mapsto x+d,\ x\mapsto-x+a_0\rangle.
\]
The translation orbits are residues modulo $d$.  Without reflections they
give one family of $d$ components, each of covering degree $m$.  With a
reflection, residues are paired by $r\mapsto a_0-r$.  If $f$ is the number
of solutions of $2r=a_0\pmod d$, there are $(d-f)/2$ paired components
of degree $2m$ and $f\le2$ exceptional fixed-residue components of degree
$m$.  This gives at most three family records per connected sewn base,
irrespective of $W$.  The same equivariant translations as in the earlier
cover module classify unmarked cover types; the two exceptional types
coincide exactly when $m$ is odd.

Orientation uses the augmented action in the affine group times
$\mathbb Z/2$.  To translation generators attach their actual character
bits; to a difference of two reflections attach the sum of their bits.
An extended-gcd calculation obtains a pair $(d,\tau)$ from these generators
and $(W,0)$.  The translation character is well defined exactly when
\[
 m\tau=0\pmod2,
 \qquad v_j=(u_j/d)\tau\pmod2
 \quad\hbox{for every translation increment }(u_j,v_j).
\]
Failure produces an orientation-reversing identity translation, so every
component is nonorientable.  Otherwise the paired components have no
reflection stabilizer and are orientable.  A fixed-residue component at
$r$ is orientable exactly when
\[
 w(R)+\frac{2r-a_0}{d}\tau=0\pmod2.
\]
These statements remain valid for the nonfaithful formal descriptions at
$W=1,2$: translation relations and actual fixed-reflection stabilizers are
both checked.

For each remaining boundary, the existing exact cycle formulas apply to
its root permutation.  A translation $x+c$ has
$q=\gcd(m,c/d)$ cycles of degree $m/q$ on a single-residue component and
$2q$ such cycles on a paired component.  A boundary reflection has $m$
two-cycles on a paired component.  On a fixed-residue component, its fixed
points solve $2y=(c-2r)/d\pmod m$; the remaining points are paired.
All multiplicities are binary integers, never allocated as individual circles.

Gluing circles does not change Euler characteristic, so
\[
 \chi(F)=\sum_{v\text{ at this root}}\chi(F_v),\qquad
 \chi(C)=\deg(C)\chi(F).
\]
If $b(C)$ is the sum of the remaining boundary-lift multiplicities, then
$2-b(C)-\chi(C)$ equals $2g(C)$ in the orientable case and the crosscap
number in the nonorientable case.  Closed output simply has $b(C)=0$.
No unverified closed-surface relator is introduced: the surface already
exists by the checked boundary-gluing construction.  The implementation
checks all resulting integrality, nonnegativity, and total-degree identities.

\begin{theorem}[Closure under full-boundary affine assembly]
Given the checked input above, every connected sewn base has at most three
binary component-family records.  The algorithm computes the degree,
orientation, genus, Euler characteristic, every remaining peripheral
degree profile, and exact component labels without expanding a sheet,
lifted boundary circle, or covering component.  It permits closed output
and preserves the original ports, seam identifications, and ordered fibre
point coordinates.
\end{theorem}
\begin{proof}
The seam-realization argument gives an actual covering of an actual
surface.  The root-action proposition computes its monodromy and orientation
character.  Dihedral orbit arithmetic gives the complete component list;
the peripheral cycles and covering Euler characteristic give its topology.
All formulas use only binary arithmetic and finitely many input records.
Provenance and point preservation are established by the explicit root and
seam transports described next.
\end{proof}

\subsection{Ports, marked points, and repeated attachment}

For a local point $(v,x)$, its root fibre coordinate is $P_v^{-1}(x)$.
The component label is its root identifier together with the canonical
dihedral orbit key.  This avoids merging components just because their
genera or unmarked types happen to agree.

For a boundary translation $x+c$, a literal lifted circle is labeled by
$x\bmod\gcd(W,c)$; for a reflection it is labeled by its one- or two-point
orbit.  Remaining-boundary keys retain the original port.  For an internal
seam, the implementation transports a point on the right side back by
$A_e^{-1}$ and uses the left boundary-cycle key together with the seam
identifier.  Thus queries on the two sides agree exactly after the
prescribed transport.  The API distinguishes an unglued boundary from an
internal seam circle: consumed boundaries cannot be queried as if they
remained on the boundary of the assembled surface.

An ordered marking consists of points $(v_i,x_i)$ in one covering component.
Transform each to the root fibre and use the complete rooted dihedral
signature from the previous cover work.  Also retain the ordered list of
piece labels $v_i$.  This additional data is necessary because a covering
isomorphism over the identity of the base cannot send a point over one
patch basepoint to a different basepoint.  The signature includes the
normalized full input model and root, so comparisons across different
assembled presentations make no implicit change-of-base claim.

A map between transitive root orbits is determined by the image of one
root point.  The existing translation/fixed/paired-orbit formulas decide
whether that image is possible.  Conjugating the resulting map by $P_v$
evaluates it at any requested point over piece $v$.  The public rooted
transport API requires its two anchors to lie over the same basepoint;
it can then transport a query point over any patch of the source component.

\texttt{with\_seams} returns a new checked assembly while retaining all
piece indices, original boundaries, and sheet coordinates.  The earlier
index remains usable.  This method recomputes the complete assembly and
does not claim incremental or amortized optimality.  In particular, marked
points survive literally as $(v,x)$ even if a new seam merges their formerly
distinct components; newly requested signatures are computed in the new
base presentation.  No unmarked type quotient is substituted for the
attaching data.

\subsection{A family showing why the seam shift must be retained}

\begin{proposition}[Identical patch topology, different closed assemblies]
Let $W=2m$ with $m\ge1$.  Take the annulus cover with core monodromy
$g(x)=x+2$ and glue its two base boundaries with $\epsilon=+1$ and
$A_b(x)=-x+b$.  For even $b$ the resulting cover consists of two Klein
bottles of degree $m$.  For odd $b$ it is one torus of degree $2m$.
\end{proposition}
\begin{proof}
The two annulus boundary monodromies are $g$ and $g^{-1}$, and
$A_b g=g^{-1}A_b$, so the seam is valid for every $b$.  The local cover
is always the same pair of annuli, each of degree $m$.  In the sewn base
the core loop has orientation bit zero and the seam loop has bit one.
If $b$ is even, reflection preserves each parity class; each is a connected
degree-$m$ cover containing an orientation-reversing stabilizer, hence is
a closed nonorientable surface with Euler characteristic zero, a Klein
bottle.  If $b$ is odd, reflection exchanges the parity classes and the
cover is connected.  An orientation-reversing group word contains an odd
number of reflections and exchanges parity, so it fixes no sheet.  The
cover is therefore orientable, closed, and has Euler characteristic zero,
hence is a torus.
\end{proof}

The proposition includes $W=2$, where the core translation is the identity
and the formal reflection may itself be the identity permutation.  These
examples are abstract surface-cover assemblies used to test topology and
attachment bookkeeping.  No embedding of the closed Klein-bottle outputs
in a knot exterior is asserted.  Their role is to prove that identical
local component topology and identical local boundary-lift profiles cannot
replace the seam transport data.  A second closed family caps a reflected
M\"obius-band cover with disk covers, giving real-projective-plane components
at fixed sheets and sphere components on paired sheets.  Two trivially
covered pairs of pants glued along all three boundaries give an orientable
genus-two base; a degree-$d$ component then has genus $d+1$.

\subsection{Bit complexity and the unresolved search problem}

Let $s$ be the number of patches, seams, generator records, and boundary
records together, and let $B$ bound the binary lengths of all input integers
and $\log(s+2)$.  Canonical peripheral words have total length $O(s)$.
Normalization, seam checks, gauge construction, and all affine changes of
base use $O(s)$ arithmetic operations on $O(B)$-bit integers, apart from
sorting piece labels, which is absorbed in the following bound.  Gcd and
extended-gcd computations give the conservative elementary bit bound
\[
 O(sB^3)
\]
for full assembly and classification.  Component degrees are at most $W$;
Euler data and genus have $O(B)$ bits because their extra factor is at most
linear in $s$.  The full binary output, including the normalized input and
gauge provenance, has $O(sB)$ bits.  There is no loop whose bound is $W$,
the number of covering components, or a boundary-lift multiplicity.

After preparation, a point's component, a seam transport, or a rooted
isomorphism evaluation uses $O(B^2)$ conservative bit work.  A port-cycle
query additionally computes a gcd and uses $O(B^3)$.  A $q$-point marked
signature uses $O(qB^2)$ work and $O(qB)$ additional output; copying,
serializing, or comparing the full model header must also account for its
$O(sB)$ size.  Independent gauge replay revalidates the patch/seam input
and tree transport equations in $O(sB^2)$ bit work.  Repeated assembly
states cost the sum of their individual bounds, not the bound of their
largest state alone.

The ordinary recognizer does not call this kernel.  Integration with a
normal-surface hierarchy still requires checked extraction of the patch
model and common fibre coordinate, representation of all relevant partial
attachments or proof that full-boundary attachments suffice, and control
of the number and size of successive repair states.  General normal-surface
pairings can be piecewise affine and need not share this global dihedral
promise.  Even one annulus cover has $W$ inequivalent ordered two-point
markings, so a constant unmarked family count does not bound marked search
states by a polynomial in $\log W$.  Polynomial bit complexity of this
assembly operation supplies one compositional primitive; it does not
establish general quasi-polynomial unknot recognition.

\subsection{Independent verification and measured scope}

The expanded oracle imports neither the gauge construction nor the gcd
classifier, peripheral-word evaluator, or marking formulas.  It constructs
literal local permutation arrays, evaluates peripherals point by point,
and forms the graph on all $(v,x)$.  Graph search gives its components,
parity propagation gives orientability, direct permutation cycles give
boundary profiles, and addition of patch Euler characteristics gives genus.
A separate labelled-edge propagation algorithm checks rooted covering
isomorphisms on that graph, preserving every local-generator and seam label.

The focused suite checks all valid self-sewn annulus actions through six
sheets, all constructed sewn-pants actions through five sheets, disk and
M\"obius cap families, random mixed orientable and nonorientable patch
assemblies, all requested port and component keys, rooted transport,
ordered repeated marks, gauge relabelings, and reversed seam descriptions.
The 14 focused tests pass.  Their recorded totals are 1547 independent expanded topology comparisons and
6216 ordered marking comparisons.  Additional tests use $W=2^{24000}$,
64-patch gauges, huge closed genus, alternative verified tree roots,
corrupt certificates, immutable seam continuation, exact hexadecimal JSON,
invalid seams, and cooperative cancellation.

The benchmark driver keeps input validation, all assembly arithmetic,
classification, and defensive summary creation in the compressed timer.
Its expanded arm includes literal graph construction and full topology.
Input generation, equality checks, destruction of returned records, and
JSON serialization are outside the timers; serialized sizes are recorded
separately.  It uses one excluded warmup, five shuffled paired measured
rounds, and an identical compressed A/A control.  The compressed topology
arm batches 100 queries per sample; the reported statistic is the median
per-query time across batches.  The expanded arm uses one complete expanded
query per sample.  The A/A topology ratios range from 0.916 to 1.016.
A second experiment
varies $W=2^{64},2^{1024},2^{4096},2^{16384}$ and uses one, eight, or 64
patches; it reports full preparation, separate gauge replay, and prepared
32-point marking queries.  Timing artifacts retain all samples, orders,
inputs or exact construction, Python/platform metadata, and source hashes.
The measured comparison is with explicit sheet expansion, not a competing
compressed orbit algorithm, and no timing in this experiment is a knot
recognition measurement.
